%!TEX root = ../main.tex
\section{Discussion and Limitations}
\label{sec:disc}

\subsection{Devel-vs-hidden-test caveat --- partition-variance equivalence framing}

All numbers in Table~\ref{tab:stack} are on the canonical \texttt{devel\_test} under \texttt{StratifiedGroupKFold}-keyed speaker-disjoint sub-splitting at SPLIT\_SEED=42. The 2017 ComParE Cold sub-challenge's $0.710$ baseline was on the hidden test set, with unspecified speaker-disjointness properties. We do not have access to the hidden test labels, so we cannot compare apples-to-apples directly.

The §\ref{ssec:shadow_robustness} shadow-split robustness analysis quantifies the canonical-split-specific lift: across 10 alternative \texttt{(devel\_val, devel\_test)} partitions of the devel split, the K=2-only mean-logit ensemble UAR varies in $[0.6443, 0.7042]$ (mean $0.6882$, $\sigma_{\text{shadow}} = 0.0169$); the canonical $0.7090$ sits at $z {=} +1.24\sigma$ above this distribution. Multi-K ensemble (§\ref{ssec:multi_k}) lifts the shadow distribution to mean $0.6940 \pm 0.0157$ ($+0.006$ paired, $9/10$ shadow splits positive) and the canonical to $0.7111$. Both our canonical numbers (K=2-only $0.7090$ or multi-K $0.7111$) and the 2017 hidden-test baseline $0.7100$ are single-partition point estimates; if the 2017 baseline has comparable partition variance ($\sigma_{2017} \sim 0.017$ is plausible given similar corpus and protocol), it could plausibly land anywhere in $[0.68, 0.74]$ on a different partition.

\textbf{We adopt the partition-variance equivalence framing as the canonical paper headline:} \emph{multi-K canonical $0.7111$ exceeds the 2017 baseline $0.7100$ by $+0.001$, and multi-K shadow-mean $0.6940$ trails the baseline by $-0.016$ ($\sim 1\sigma_{\text{shadow}}$); both gaps are smaller than the partition variance of either single-partition estimate. We consider the systems statistically equivalent within partition variance, with our locked configuration in the slight-canonical-favourable / slight-shadow-unfavourable arm of that equivalence.} This framing is stronger and more defensible than the previously-listed conservative / standard / generous framings because it (a) acknowledges and quantifies the canonical favourability ($+1.24\sigma$ on K=2-only, $+1.09\sigma$ on multi-K, $\Delta\,+0.017$ to $+0.021$ over the shadow mean); (b) attributes the comparison's uncertainty to the right source (single-partition noise, not measurement error); (c) preempts the obvious reviewer question ``did you overfit devel?'' by reporting the shadow distribution alongside the canonical headline.

\textbf{Same-protocol classical-baseline-architecture replication.} Section~\ref{ssec:multi_k} reports a re-run of the 2017 ComParE Cold sub-challenge baseline architecture (6373-d ComParE-2016 functionals + regularised cold-LR cold probe) under our speaker-grouped subsplit. Standalone shadow-mean UAR $= 0.5853 \pm 0.008$, meaningfully below both FM-A2.5 ($\sim 0.656$) and the 2017 hidden-test baseline ($0.7100$). The $0.5853 \to 0.7100$ gap of $\Delta\,+0.125$ is the joint effect of the speaker-grouped subsplit being stricter than the 2017 hidden-test split protocol (cross-speaker leakage closes) and the SVM-vs-LR architectural choice; the protocol-stricter split likely dominates because SVM-vs-LR architectural differences rarely produce $+0.10+$ UAR gaps on this scale of data. \emph{We therefore quantify the FM-based system's lift over the classical baseline architecture on the same protocol as $\Delta\,+0.109$ shadow-mean ($0.6940$ vs $0.5853$).} This same-protocol comparison provides a stronger and more defensible C1 contribution claim than the apples-to-oranges canonical-vs-baseline comparison: the FM-based pipeline dominates the classical-feature-family architecture by $\sim 11$\,pp UAR on the same strict-protocol evaluation, independent of the canonical-vs-baseline single-partition equivalence framing.

The multi-K-with-K=3 (adding ComParE-LR as a fusion candidate) reaches canonical UAR $0.7126$, the highest single-partition canonical achieved across the entire project ($\Delta\,+0.0026$ vs the 2017 baseline). Shadow-mean is unchanged ($0.6941$, $\Delta\,+0.00002$, positive on $5/10$ shadow splits = coin-flip), so this configuration is paper-supplementary rather than headline. Future work with hidden-test access could resolve the canonical-vs-baseline comparison directly; until then, the shadow distribution + the same-protocol classical-baseline replication are the right uncertainty + dominance frameworks for the headline claim.

\subsection{Exploratory-vs-confirmatory framing}

The K=2 G\_other candidate selection (Section~\ref{ssec:k2}) tested 5 candidates $\{G1, G2, G3, G5, G6\}$ on devel\_test and locked G5\_modulation as the winner. Strict reading: this is a ``best-of-5'' devel-test selection. The 5-seed expansion (Section~\ref{ssec:k2_5seed}) at \emph{fixed} G5 on 2 unseen-during-selection seeds $\{999, 31337\}$ is confirmatory; the new seeds' K=2 lifts ($+0.0092$, $+0.0095$) sit within the original 3-seed range ($+0.0075$ to $+0.0137$), bounding the selection effect. We are explicit about this framing to avoid the multiple-comparison-on-devel critique. The exhaustive K=2 sweep across all 6 admitted A5a groups + the K=3 with full eGeMAPSv02 confirms G5 as the unique winner under the locked-$\beta$ methodology.

\subsection{Ensemble caveats}

The K=2 5-seed mean-logit ensemble is a single-number evaluation, not $\sigma$-bearing. The $\sigma$-bearing canonical headline stays the per-seed K=2 single mean (0.7037 $\pm$ 0.0060). The ensemble is a paper-supplementary ablation row showing $+0.005$ UAR from averaging at 5$\times$ inference cost. Per-seed single is the deployment-relevant number; the ensemble is the upper bound on this stack.

\subsection{Methodological honesty about lr $\times$ init grid}

The A2.5 lr stress test produced a clarification of the original ``flat landscape'' claim: the layer-weight subspace has gradient signal at lr$\times 10$ ($\cos(\text{init}, \text{final})$ drops to $0.83$ at lr$\times 10$, $0.21$ at lr$\times 100$). The lr $\times$ init grid (uniform vs honest-prior $\times$ lr$\times \{1, 3, 5, 10\}$, 24 retrains, 3 seeds per cell) shows the honesty prior is a \emph{distinct attractor} that uniform initialisation cannot reach via lr alone --- uniform plateaus at UAR $0.6447$--$0.6466$ across all tested lr while honesty-prior reaches $0.657$--$0.664$ at every lr. We report A2.5-default (honest-prior + lr$\times 1$) as canonical because higher-lr trades $+0.007$ UAR for $+0.016$ LR-probe leak --- a Pareto-bad trade for the de-confounding paper narrative. The grid serves as ablation evidence; the per-cell numbers and the cosine table across the lr-axis are in Appendix~\ref{app:lr_init_grid}.

\subsection{Limitations}

\textit{(L1) Speaker-probe noise floor at this corpus scale (M14).} URTIC at $k{=}210$ pseudo-speakers with $\sim$40 chunks/speaker produces a speaker-probe top-1 ceiling around 0.085 even with the strongest tractable discriminator. ``Probe top-1 dropped from $X$ to $Y$'' framings common in the de-confounding literature are potentially misleading on similar-scale corpora; we report disc-ceiling + memo gap alongside every probe number we cite. Future work on larger-corpus or fine-grained-speaker variants could quantify the tightening of the noise floor.

\textit{(L2) Frozen-backbone architectural commitment.} We test no full-transformer fine-tune. The de-confounding ladder closure is therefore ``at every frozen-backbone scope tractable in budget,'' not ``at every architectural level full stop.'' A full fine-tune may give a different verdict for both contrastive (A6) and adversarial (A7) interventions. Documenting this scope limit is part of the project's epistemic honesty; we do not extrapolate beyond it. Similarly, the A5b K=3 multi-FM test (Section~\ref{ssec:k3_hubert}) used HuBERT-base mean-pooled — the simplest and cheapest variant; deeper multi-FM tests (HuBERT-base \emph{with} learned layer-weighted softmax, or HuBERT-large) remain future work and could in principle admit at K=3 even though the simplest variant cleanly fails per M14.

\textit{(L3) Pseudo-speaker labels (k=210 KMeans on ECAPA).} We use ECAPA-TDNN embeddings clustered into $k{=}210$ as the speaker-grouping proxy. URTIC's true speaker count is $\sim$630; our pseudo-speakers approximate this. A finer or coarser $k$ may shift probe ceilings; the M14 finding is that the qualitative pattern (memorisation dominates) is stable across $k$.

\textit{(L4) Devel\_val sub-split for early stop is also speaker-disjoint from devel\_test, but not from training in the sense of having any cross-validation discipline.} The early-stop split was fixed at \texttt{val\_frac=0.50} with seed 42; a different seed/frac may slightly shift the locked-$\beta^*$ choices. The stability of $\beta^* \in \{6, 8\}$ across all 5 seeds for K=2 suggests this is a minor concern.

\textit{(L5) Class-balanced sampling vs natural class distribution.} A2 / A2.5 / K=1 / K=2 all use balanced sampling. The 9.5\,\% cold rate at deployment time may produce different operating-point trade-offs; the recall pattern shift in the ensemble (Section~\ref{ssec:ensemble}) suggests the K=2 ensemble is the most cold-recovery-friendly variant.
